\section{Evaluación de la Solución}
\subsection{Puntos fuertes y débiles}
Entre los puntos fuertes de la solución se destacan la separación entre la recepción de los eventos y su procesamiento, que permite tolerar la indisponibilidad del componente central y de los nodos periféricos sin perder información; el control de acceso aplicado en la capa de servicios y no solo en las interfaces; la posibilidad de seguir cada operación de punta a punta mediante trazas distribuidas; y un entorno de desarrollo completo que se levanta con un solo comando.

Entre los puntos débiles, al momento de esta entrega quedan pendientes la regla de desvío de ruta, que requiere calcular el recorrido esperado con soporte geoespacial; la consulta a la PDI en el primer ingreso; la asignación de choferes a una empresa; el filtro por ubicación de los reportes; la validación de viajes solapados de un mismo vehículo; y la notificación de los casos por correo. Además, la performance se evaluó solo en forma preliminar.

\subsection{Volumen de tracking y performance de la ingesta}
El volumen objetivo se estimó a partir de los 17.869 vehículos con motor registrados para el transporte profesional de carga \cite{BID2022}, de las 2.050 horas de trabajo por año que el MTOP considera para un camión pesado \cite{INALOG2021} y de un factor de hora pico de entre el 10~\% y el 16~\% del tránsito diario \cite{PennDOT2021}: en la hora pico habría entre 10.000 y 16.000 vehículos en viaje. Se tomaron 15.000 vehículos que reportan una posición cada 30 segundos, una frecuencia que fija la plataforma: 500 posiciones por segundo, con picos del triple.

Una prueba preliminar, con todos los componentes y el generador de carga en una misma computadora de 8 núcleos, dio los resultados de la Tabla~\ref{tab:prueba-preliminar}. Hasta 250 posiciones por segundo la cola se mantiene acotada; con 500, el central procesa en promedio lo que llega, pero la computadora se satura. El procesamiento de cada posición no supera los 8~ms en el percentil 95, por lo que el cuello de botella, identificado con las métricas de la plataforma de observabilidad, es la CPU compartida y no la ingesta. Queda pendiente repetir la prueba con el generador separado y el central en la nube, medir el frontoffice bajo carga y verificar el escalado horizontal.

\begin{table}
\caption{Prueba preliminar de la ingesta de tracking.}\label{tab:prueba-preliminar}
\centering
\begin{tabular}{|r|r|r|r|}
\hline
Objetivo (pos/s) & Procesadas (pos/s) & Cola máxima & Procesamiento p95 \\
\hline
100 & 88 & 1.002 & 38 ms \\
\hline
250 & 279 & 2.192 & 8,4 ms \\
\hline
500 & 521 & 9.932 & 7,7 ms \\
\hline
\end{tabular}
\end{table}

\subsection{Resiliencia ante fallas}
Con el componente central detenido, el nodo de tracking siguió publicando y las posiciones esperaron en la cola; al reiniciarlo, se procesaron todas. Con PostgreSQL detenido durante un minuto, el central se recuperó sin intervención y ninguna posición llegó a la cola de mensajes muertos. Al detener un nodo de balanzas, el caso de exceso de carga se registró con los pesajes del otro nodo, y el pesaje faltante se agregó al mismo caso cuando el nodo volvió.
